\documentclass{article}
\usepackage{mathnotes}

\title{Properties of Vectors}
\description{Fundamental vector operations in Rⁿ including magnitude, addition, scalar multiplication, and linear combinations with definitions of span and vector spaces.}

\begin{document}

\section{Vector Space}

\begin{definition}[Vector Space]
A \textbf{vector space} over a \dref{field} $F$ is a \dref{set} $V$ together with two operations:

\begin{itemize}
\item Vector addition: $V \times V \to V$
\item Scalar multiplication: $F \times V \to V$
\end{itemize}

satisfying the following axioms:

\begin{enumerate}
\item $(V, +)$ is an \dref{abelian} \dref{group}
\end{enumerate}

\begin{enumerate}
\item Scalar multiplication is associative: $a(bv) = (ab)v$
\end{enumerate}

\begin{enumerate}
\item Distributive laws hold
\end{enumerate}

\begin{enumerate}
\item Identity: $1v = v$ for all $v \in V$
\end{enumerate}
\end{definition}

\section{Scalars}
\begin{definition}[Scalar]
A \textbf{scalar} is an \dref{element} of a \dref{field} used to define a \dref{vector-space}.
\end{definition}

\begin{note}
Typically, especially in physics, a \dref{scalar} is simply a number, especially a \dref[real number]{real-numbers}.
\end{note}

\section{Vectors}
\begin{definition}[Vector]
A \textbf{vector} is an \dref{element} in a \dref{vector-space}.
\end{definition}

\begin{note}
Typically, in physics and many other applications, when we say \dref{vector} we mean a \dref{vector} in $\mathbb{R}^n,$ which is an ordered tuple of real numbers, $\vec{x} = (x_1, x_2, \cdots, x_n).$ This is the algebraic representation of vector $\vec{x}$.

The geometric representation of vector $\vec{x}$ is, in an $n$ dimensional space, an arrow or directed line segment. When starting from the origin, it would be a line segment from $\vec{0}$ to the point $(x_1, x_2, \cdots, x_n)$.

Physically, a vector has a magnitude and a direction.
\end{note}

\begin{definition}[Free vector]
A \dref{vector} whose endpoints are not fixed at particular \dref{points} is called a \textbf{free vector.}
\end{definition}

\begin{definition}[Bound vector]
A \dref{vector} with a fixed endpoint is called a \textbf{bound vector.}
\end{definition}

\section{Properties of Vectors}

\begin{definition}[Magnitude]\synonyms{norm, length}
Let $\vec{x} = (x_1, x_2, \cdots, x_n) \in \mathbb{R}^n$.  The \textbf{magnitude}, or length, $\vec{x}$ is denoted as $|| \vec{x} ||$ and is defined as:

\[ ||\vec{x}|| = \sqrt{ {x_1}^2 + {x_2}^2 + \cdots + {x_n}^2 } = \sqrt{\vec{x} \cdot \vec{x}} = \sqrt{\sum_{i=1}^n x_i^2} \]
\end{definition}

\begin{note}
This is essentially the Pythagorean theorem in $n$ dimensions; in $\mathbb{R}^2$ the magnitude of a vector corresponds to the length of the hypotenuse of a right triangle whose other sides are of length $x_1$ and $x_2.$

While we define \textbf{norm} to be equivalent to magnitude here, this is actually a special case of the more general concept of a norm - there are other norms we could define on $\mathbb{R}^n,$ but I don't have need to explore that yet.
\end{note}

\begin{definition}[Direction]
The \textbf{direction} of a \dref{vector} can be specified by the angle between it and some fixed reference, such as the $x$-axis.
\end{definition}

\begin{definition}[Zero Vector]
The \textbf{zero vector} $(0, 0, \cdots, 0)$ is denoted as $\vec{0}$ and has no \dref{direction}.
\end{definition}

\begin{definition}[Vector Equality]
Two vectors are equal if they have the same coordinates (or equivalently, the same \dref{length} and \dref{direction}).
\end{definition}

\begin{definition}[Unit Vector]
If a \dref{vector} has a \dref{length} of 1, i.e. if $||x|| = 1$, we say the \dref{vector} is a \textbf{unit vector}.
\end{definition}

\begin{definition}[Vector Addition]
We perform \textbf{vector addition} by adding two \dref{vectors} $\vec{x} = (x_1, x_2, \cdots, x_n)$ and $\vec{y} = (y_1, y_2, \cdots, y_n)$ according to the following rule:

\[ \vec{x} + \vec{y} = (x_1 + y_1, x_2 + y_2, \cdots, x_n + y_n) \]

that is, by making a new vector where the coordinates are the sums of the respective coordinates in the vectors being summed.
\end{definition}

\begin{note}
Geometrically, \dref{vector-addition} connects \dref{vectors} head to tail.
\end{note}

\begin{definition}[Vector Subtraction]
Similarly, \textbf{vector subtraction} can be performed as:

\[ \vec{x} - \vec{y} = (x_1 - y_1, x_2 - y_2, \cdots, x_n - y_n) \]
\end{definition}

\subsection{Multiplication}
\begin{definition}[Vector Multiplication by a Scalar]
We can multiply \dref{vectors} by a \dref{scalar}. Given the \dref{scalar} $c$:

\[ c \vec{x} = (c x_1, c x_2, \cdots, c x_n) \]
\end{definition}

\begin{definition}[Inner product]\synonyms{"Dot Product", "Scalar product of two vectors"}
If $\vec{x}$ and $\vec{y}$ are vectors in $\mathbb{R}^n,$ then their \textbf{inner product} is defined as

\[ \vec{X} \cdot \vec{y} = \sum_{i = 1}^n x_i y_i. \]
\end{definition}

\begin{theorem}\label{cos-characterization-of-dot-product}
The \dref{dot-product} of $\vec{u}$ and $\vec{v}$ is

\[ \vec{u} \cdot \vec{v} = |\vec{u}||\vec{v}|\cos{\theta} \]

where $\theta$ is the angle between $\vec{u}$ and $\vec{v}.$
\end{theorem}

\begin{definition}[Perpendicular]\synonyms{orthogonal}
Two vectors $\vec{u}$ and $\vec{v}$ are said to be \textbf{perpendicular} or \textbf{orthogonal} if the angle between them is $\pi/2$ radians, or, equivalently, if the inner product $\vec{u} \cdot \vec{v} = 0.$
\end{definition}

\begin{definition}[Parallel]
Two vectors $\vec{x}$ and $\vec{y}$ are said to be \textbf{parallel} if $\vec{x}$ is a scalar multiple of $\vec{y}$, i.e., if there exists some scalar $c$ where $\vec{x} = c \vec{y}.$
\end{definition}

\begin{definition}[Cross Product]\synonyms{"Vector Product", "Outer Product"}
The \textbf{cross product} $\vec{a} \times \vec{b}$ (read "a cross b") of two \dref{vectors} $\vec{a}$ and $\vec{b}$ is the \dref{vector} $\vec{v}$ denoted by

\[ \vec{v} = \vec{a} \times \vec{b}. \]

I. If $\vec{a} = \vec{0}$ or $\vec{b} = \vec{0},$ then we define $\vec{v} = \vec{a} \times \vec{b} = \vec{0}.$

II. If both \dref{vectors} are nonzero \dref{vectors}, then \dref{vector} $\vec{v}$ has the \dref{length}

\[ \tag{1} |\vec{v}| = |\vec{a} \times \vec{b}| = |\vec{a}||\vec{b}|\sin{\theta}, \]

where $\theta$ is the angle between $\vec{a}$ and $\vec{b}.$ Furthermore, $\vec{a}$ and $\vec{b}$ form the sides of a parallelogram on a plane in space. The area of this parallelogram is precisely given by (1), such that the length $|\vec{v}|$ of the \dref{vector} $\vec{v}$ is equal to the area of the parallelogram.

III. If $\vec{a}$ and $\vec{b}$ lie in the same straight line, i.e. $\vec{a}$ and $\vec{b}$ have the same or opposite \dref{directions}, then $\theta$ is $0 \degree$ or $180 \degree$ so that $\sin \theta = 0.$ In that case $|\vec{v}| = 0,$  so that $\vec{v} = \vec{a} \times \vec{b} = \vec{0}.$

IV. If cases I and III do not occur, then $\vec{v}$ is a nonzero vector. The \dref{direction} of $\vec{v} = \vec{a} \times \vec{b}$ is \dref{perpendicular} to both $\vec{a}$ and $\vec{b}$ such that $\vec{a}, \vec{b}, \vec{v},$ precisely in this order, form a right-handed triple
\end{definition}

\begin{remark}
If $\vec{a} = (a_1, a_2, a_3)$ and $\vec{b} = (b_1, b_2, b_3),$ then we use the symbolic determinant to find the \dref{cross-product}, in this fashion:

\[ \vec{v} = \vec{a} \times \vec{b} =  \begin{vmatrix}\vec{i} & \vec{j} & \vec{k}\\a_1 & a_2 & a_3\\b_1 & b_2 & b_3\end{vmatrix} = \begin{vmatrix} a_2 & a_3 \\ b_2 & b_3 \end{vmatrix} \vec{i} - \begin{vmatrix} a_1 & a_3 \\ b_1 & b_3 \end{vmatrix} \vec{j} + \begin{vmatrix} a_1 & a_2 \\ b_1 & b_2 \end{vmatrix} \vec{k}. \]
\end{remark}

\includedemo{cross-product}

\begin{definition}[Linear Combination]
Let $\vec{v_1}, \vec{v_2}, \cdots, \vec{v_n} \in \mathbb{R}^n$ and $c_1, c_2, \cdots, c_3 \in \mathbb{R}.$ Then, the vector

\[ \vec{v} = c_1 \vec{v_1} +  c_2 \vec{v_2} + \cdots + c_n \vec{v_n} \]

is called a \textbf{linear combination} of $\vec{v_1}, \vec{v_2}, \cdots, \vec{v_n}.$

Let $\vec{v_1}, \vec{v_2}, \cdots, \vec{v_n} \in \mathbb{R}^n.$ The set of all linear combinations of $\vec{v_1}, \vec{v_2}, \cdots, \vec{v_n}$ is called their \textbf{span}, denoted $\Span{\left(\vec{v_1}, \vec{v_2}, \cdots, \vec{v_n}\right)}.$ 

That is:

\[ \Span{\left(\vec{v_1}, \vec{v_2}, \cdots, \vec{v_n}\right)} = \left\{ \vec{v} \in \mathbb{R}^n : \vec{v} = c_1 \vec{v_1} + c_2 \vec{v_2} + \cdots + c_n \vec{v_n} ~ \text{for some scalars} ~ c_1, c_2, \cdots, c_n \right\} \]
\end{definition}

\end{document}
